\documentclass{beamer}
%\usepackage{split}

\usepackage{amsmath}
\usepackage{amsfonts}
\usepackage{braket}
\usepackage{subcaption}
\usepackage{tikz}
\usepackage{xcolor}
\usepackage{svg}
\usetikzlibrary{positioning,shapes,arrows.meta}
%\usetikzlibrary{positioning,fit,backgrounds}
\usepackage{quantikz}
\usepackage[linesnumbered]{algorithm2e}
\usepackage{hyperref}
%\usepackage{hyperref,xcolor}
%\usepackage[ocgcolorlinks]{ocgx2}
\usepackage{cleveref}
\usepackage[backend=biber,style=alphabetic,sorting=ynt]{biblatex}
\addbibresource{./sample.bib}


\input{newcommands}


\newcommand*{\QACze}{ \mathbf{QAC}_{0} }
\newcommand*{\QNCzef}{ \mathbf{QNC}_{0,f} }
\newcommand*{\QNCze}{ \mathbf{QNC}_{0} }
\newcommand*{\QNCon}{ \mathbf{QNC}_{1} }
\newcommand*{\NCon}{ \mathbf{NC}_{1} }
\newcommand*{\noiseQNCon}{ noisy-$\QNCon$ }
\newcommand*{\QNC}{ \mathbf{QNC} }
\newcommand*{\QNCG}{ \mathbf{QNC_G} }
\newcommand*{\NC}{\mathbf{NC}}
\newcommand*{\QNCiG}{\mathbf{QNC_{G,i}}}



\usetikzlibrary{patterns, shapes.arrows}
\usepackage{tkz-graph}
\usetikzlibrary{decorations.pathmorphing}
\usepackage{tkz-berge}

\usetikzlibrary{intersections}

\usepackage{cancel}
\begin{document} 

%\newcommand*{\Tr}{\textbf{Tr }}


\begin{frame}
  \title{Quantum Toom's Contours}
    \author{David Ponarovsky}
    \date{\today}
    \titlepage
\end{frame}


\begin{frame}

\frametitle{Introduction}
\begin{block}{Today:}
\begin{itemize}
  \item Proof. 
    \begin{enumerate}
      \item Walktrough the math 
   %   \item A 
  %    \item B
    \end{enumerate}
%  \item New experimental results.  
 % \item Answering the code teleportation.  
\end{itemize}
\end{block}
\end{frame}

% ------------------------------------------------ $ 
  
\begin{frame}
  \frametitle{Our Code}

  \begin{block}{Code}

    Consider the representation of the $d$-dimensional toric by the cayley graph of the $\left( \mathbb{Z}_{n}^{d}, + \right)$ group generated by the elements $S = \{ 1_{i} : i \in [d] \}$ \footnote{where $1_{i}$ is the vector consisting of one at the $i$th coordinate and zero elsewhere}.

    In our code, bits are set on the plaquette, and the edges correspond to checks that restrict the parity of the bits on their support to be even.

  \end{block}


\end{frame}

\begin{frame}
  \frametitle{Local Decoder} 

  
  \begin{block}{Positive edge}
    We say that an edge $\{x, y\}$ is positive relative to a plaquette $\{v, g_{1}v, g_{2}g_{1}v, g_{2}v\}$, defined by a rooted vertex $v \in \mathbb{Z}_{n}^{d}$ and the generators $g_{1}\neq g_{2} \in S$, if $x = v$ and $y$ is either $g_{1}v$ or $g_{2}v$.

  \end{block}

  \begin{block}{Toom's rule}
We generalize Toom's rule in the following manner: Each plaquette asks if both of its positive edges point to a syndrome, namely unsatisfied. If so, then the bit flips itself; for convenience, we will think of the vertex at the end of the positive edges as the vertex which makes the decision and flips the plaquette.
\end{block}

\end{frame}


\begin{frame}
  \frametitle{Title} 


Let $\sigma^{(t)}$ be a domain wall at time $t$, and let $\Gamma^{t} = \{ v : \exists \{v,u\} \in \sigma^{t} \}$. We denote by $\phi^{t}$ the subsets of bits which are flipped as a result of applying the Toom's rule at time $t$, and by $\phi_{v}$ the bits which are flipped by vertex $v$. We say that $u$ is added to $\Gamma^{t+1}$ by $v$ at time $t$, if $u\in \Gamma^{t+1}$ and $u \in V(\phi_{v})$. When it clear form the context, we will omit the time, and say that $u$ is added by $v$. 


\end{frame}





\end{document}
